\subsection{A smooth eight-cubic seed at relaxed length}
\label{app:eight-cubic-smooth-seed}

\begin{theorem}
\label{thm:eight-cubic-smooth-seed}
Over a number field, and over arbitrarily large prime fields, there is a
word on $16$ distinct coordinates whose maximum agreement with cubics is
$7$ and whose complete nearest list consists of exactly eight cubics.
Consequently, over arbitrarily large prime fields there is a length-$32$,
dimension-$7$ Reed--Solomon word whose maximum agreement is $14$ and
whose complete nearest list consists of exactly eight codewords. This agreement exceeds the audited
first-order curve at rate $7/32$. The construction has fixed length and
list size; no growing-list or asymptotic-tightness assertion follows.
\end{theorem}

\begin{proof}
We first give a smooth point of the incidence equations modulo $17$.
The coordinates are $x_j=j$, $1\le j\le16$, and the received word is
\[
 (9,3,12,15,0,14,8,13,11,13,11,8,10,14,10,13).
\]
The eight cubics $P_i(X)=\sum_{r=0}^3 c_{ir}X^r$ and their exact agreement
sets are as follows; entries and evaluations in this table are modulo $17$.
\[
\begin{array}{c|rrrr|l}
 i&c_{i0}&c_{i1}&c_{i2}&c_{i3}&\text{agreement coordinates}\\\hline
 1&1&10&7&2&3,5,6,10,14,15,16\\
 2&16&11&12&4&1,4,8,11,13,15,16\\
 3&13&4&7&5&5,7,9,11,13,14,15\\
 4&3&1&15&7&1,4,5,10,11,12,14\\
 5&0&10&2&10&3,4,6,9,11,12,15\\
 6&15&12&5&12&3,5,8,9,12,13,16\\
 7&14&15&0&14&1,2,3,7,11,12,13\\
 8&16&1&13&15&2,4,7,9,12,14,16
\end{array}
\]
All nodes are distinct and nonzero, all leading coefficients are nonzero
and distinct, and each displayed polynomial agrees exactly seven times.
For completeness, interpolate every four-point received subset. The
$\binom{16}{4}=1820$ subsets produce $1022$ distinct cubics, whose maximum
agreement is seven, with exactly the eight maximizers displayed above.
This exhausts every possible cubic with at least four agreements, so it
proves completeness even over the algebraic closure of $\mathbb F_{17}$.
An independent verifier uses four-by-four Vandermonde elimination rather
than the search algorithm.

Introduce $64$ variables: the $32$ coefficients $c_{ir}$, the $16$
coordinates $x_j$, and the $16$ received values $w_j$. Impose the $56$
equations
\[
 \sum_{r=0}^3c_{ir}x_j^r-w_j=0\qquad(j\text{ in the displayed support of }i).
\]
A row of their Jacobian has entries $1,x_j,x_j^2,x_j^3$ in the four
coefficient columns of $P_i$, entry $P_i'(x_j)$ in its coordinate column,
and $-1$ in its received-value column. Select all $32$ coefficient
columns, the first $14$ coordinate columns, and the first $10$
received-value columns. The resulting $56$-by-$56$ determinant is
\[
                         14\pmod {17}.
\]
This determinant was independently reconstructed from the table.

Fix the remaining eight variables, namely $x_{15},x_{16}$ and
$w_{11},\ldots,w_{16}$, to their displayed integer representatives.
The selected determinant is a unit, so multivariate Hensel lifting gives
a solution in $\mathbb Z_{17}$ to the resulting square system, with
invertible Jacobian. The solution lies on a zero-dimensional component
of the characteristic-zero system: its local dimension is zero by that
Jacobian. Its coordinates are consequently algebraic over $\mathbb Q$.
They generate a number field embedded in $\mathbb Q_{17}$, equipped with
the indicated reduction. Distinct nodes, nonzero nodes, distinct cubics,
and all unselected disagreement inequalities persist, since their
reductions are nonzero.

The lifted word still has the same complete nearest list. Indeed, any
cubic with at least seven agreements is determined by four of them.
The corresponding interpolation denominators are units at the chosen
place, so the cubic is integral there and reduces to a cubic with at
least seven agreements in the table's word. It must reduce to one of
the eight listed cubics, and its seven matching indices must be exactly
that cubic's seven-index support. Both characteristic-zero polynomials
agree with the word on this support, so four-point uniqueness makes
them identical. Agreement greater than seven would contradict the
modular maximum. This also proves the assertion over an algebraic
closure of the number field, since four-point interpolation already
puts every such cubic in the number field.

Now adjoin the square roots of the sixteen lifted coordinates. They
are nonzero and pairwise distinct, so the $32$ points
$\{\pm\sqrt{x_j}:1\le j\le16\}$ are distinct. Give both points above
$x_j$ the received value $w_j$. The eight distinct polynomials
$P_i(T^2)$ have degree six and exactly fourteen agreements each.
We verify its completeness first over $\overline{\mathbb F}_{17}$.
A non-even competitor has the unique form
$Q(T)=E(Y)+T O(Y)$, $Y=T^2$, with $\deg E\le3$, $\deg O\le2$ and
$O\ne0$. At most two fibers can match at both signs, so fourteen
agreements require at least twelve distinct matched fibers. At each,
\[
 w_j^2+B(x_j)w_j+C(x_j)=0,\qquad
 B=-2E,\quad C=E^2-YO^2.
\]
There are eleven coefficients, four in $B$ and seven in $C$.
Enumerating the $1820$ twelve-fiber subsets gives $1686$ inconsistent
rank-eleven systems, $132$ consistent rank-eleven systems, and two
inconsistent rank-ten systems. The consistent systems give $84$
distinct solutions, all defined over $\mathbb F_{17}$ even when the
candidate coefficients are allowed in its algebraic closure.
Every solution violates at least one of
\[
 C(0)=E(0)^2,\qquad [Y^6]C=([Y^3]E)^2,
\]
which are necessary since $\deg(YO^2)\le5$. Thus there is no such
non-even competitor. Even competitors are exactly the eight supplied
by the complete cubic source. This proves maximum agreement fourteen
and complete nearest-list size eight over $\overline{\mathbb F}_{17}$.

For characteristic-zero completeness, use the unramified quadratic
extension of $\mathbb Q_{17}$. All sixteen nonzero residue coordinates
have square roots in its residue field $\mathbb F_{17^2}$; their two
roots are simple and distinct. The $32$ lifted nodes are consequently
unit-separated. A degree-six competitor with fourteen agreements has
integral coefficients by interpolation on seven matching nodes. Its
reduction must be one of the eight just classified, with exactly that
codeword's fourteen-index support. Seven-point uniqueness then identifies
the characteristic-zero competitor with the corresponding $P_i(T^2)$.
This excludes also a non-even polynomial whose odd part would reduce
to zero.

Take a finite Galois number-field extension containing this construction.
Outside finitely many primes, all coefficients and points reduce, and
all displayed distinctness and disagreement conditions persist.
For completeness one likewise excludes the finitely many nonzero
evaluation differences arising from four-point interpolation on the
source and seven-point interpolation on the pullback.
There are arbitrarily large rational primes splitting completely in
this extension; their residue fields are prime fields and contain all
$32$ pullback coordinates. Thus the construction and both complete nearest-list assertions
specialize as claimed.

Finally, for $\rho=7/32$ and $a=14/32$, one has
$\rho>11-3\sqrt{13}$ and
\[
 (8-\rho)a^2-6\rho a+\rho(4\rho-5)=\frac{105}{8192}>0.
\]
The audited high-rate first-order curve is the positive root of this
quadratic, so the agreement is strictly above it. The capacity gap is
$a-\rho=7/32$. These are finite parameters, not an unbounded-list family.
\end{proof}

The reproducibility directory is
\path{research/relaxed_eight_cubic_route/}.
The seed, supports, and selected Jacobian columns are in
\path{local_search.json}; the independent incidence, determinant, and
complete four-point census replay is
\path{independent_verify.py} with receipt
\path{independent_verify.json}. The independent replay uses only exact
modular arithmetic and verifies the first-order margin as a rational
number.

The additional norm census is \path{quadratic_pullback_gate.json}.
Its independent replay, \path{verify_quadratic_pullback.py}, reverses the
coefficient-column order and uses forward elimination and back substitution;
\path{quadratic_pullback_gate.verified.json} records all $84$ coefficient
obstructions. No square-root test or numerical approximation is needed.
